\documentclass[11pt]{report}

\usepackage[a4paper,margin=1in]{geometry}
\usepackage{amsmath,amssymb,amsthm}
\usepackage{enumitem}
\usepackage{url}
\usepackage[colorlinks=true,linkcolor=blue,citecolor=blue,urlcolor=blue]{hyperref}

\title{Personal Notes on Jordan Quillen (Co)homology}
\author{}
\date{\today}

\theoremstyle{definition}
\newtheorem{definition}{Definition}[chapter]
\newtheorem{example}{Example}[chapter]
\newtheorem{note}{Note}[chapter]
\newtheorem{question}{Question}[chapter]

\theoremstyle{plain}
\newtheorem{theorem}{Theorem}[chapter]
\newtheorem{lemma}{Lemma}[chapter]

\begin{document}
\maketitle
\tableofcontents

\chapter{Background}

\section{Purpose}

This chapter is a private learning layer for the mathematical background behind
intrinsic Quillen homology and cohomology of Jordan algebras. It is not a claim
registry and should not be treated as proof evidence.

\begin{note}
  Let \(k\) be a field of characteristic \(0\), and let
  \(\operatorname{Jord}_k\) denote a suitable category of Jordan algebras over
  \(k\), for example nonunital Jordan algebras, simplicial Jordan algebras, or
  dg Jordan algebras over the Jordan operad.

  The word \emph{intrinsic} means that the construction is carried out inside the
  category of Jordan algebras itself. In particular, one does not first pass from
  a Jordan algebra \(J\) to an associated Lie algebra, such as the Tits--Kantor--
  Koecher Lie algebra, and then take Lie algebra homology.

  In the nonunital absolute setting, the relevant abelianization functor is the
  functor of indecomposables
  \[
      Q(J) = J/J^2,
  \]
  where
  \[
      J^2 := \operatorname{span}_k\{xy : x,y\in J\}.
  \]
  The quotient \(J/J^2\) has zero multiplication, and it is the universal
  zero-multiplication Jordan algebra receiving a morphism from \(J\).

  To define Quillen homology, one replaces \(J\) by a cofibrant replacement
  \[
      \widetilde{J} \xrightarrow{\sim} J
  \]
  in the chosen homotopical category of Jordan algebras, and then applies \(Q\).
  Thus the derived indecomposables are
  \[
      \mathbb{L}Q(J) := Q(\widetilde{J})
      = \widetilde{J}/(\widetilde{J})^2.
  \]
  The intrinsic Jordan Quillen homology of \(J\) is
  \[
      H^{\operatorname{Jord}}_n(J)
      :=
      H_n\bigl(\mathbb{L}Q(J)\bigr).
  \]

  More generally, for a morphism of Jordan algebras
  \[
      A \longrightarrow B,
  \]
  one expects a relative cotangent complex
  \[
      \mathbb{L}^{\operatorname{Jord}}_{B/A}.
  \]
  If \(M\) is an appropriate \(B\)-module, for example a Beck module or Jordan
  module, then the intrinsic Jordan Quillen cohomology is defined by
  \[
      D^n_{\operatorname{Jord}}(B/A,M)
      :=
      H^n\operatorname{Hom}_B
      \bigl(
          \mathbb{L}^{\operatorname{Jord}}_{B/A},
          M
      \bigr).
  \]
  Dually, one may define homology with coefficients by
  \[
      D^{\operatorname{Jord}}_n(B/A,M)
      :=
      H_n\bigl(
          \mathbb{L}^{\operatorname{Jord}}_{B/A}
          \otimes_B M
      \bigr),
  \]
  or, more carefully, by using the appropriate derived tensor product.

  The low-degree interpretation should be analogous to the commutative and Lie
  cases:
  \[
      D^0_{\operatorname{Jord}}(B/A,M)
      \cong
      \operatorname{Der}_A^{\operatorname{Jord}}(B,M),
  \]
  while
  \[
      D^1_{\operatorname{Jord}}(B/A,M)
  \]
  should classify square-zero extensions of \(B\) by \(M\), subject to the precise
  choice of the coefficient category.
\end{note}

\section{Working Setting}

The default ground field is a field \(k\) of characteristic \(0\). The current
working category is usually nonunital Jordan algebras, with augmented unital
algebras treated through their augmentation ideals when that convention is
appropriate.

\begin{definition}
A Jordan algebra is a commutative algebra whose product satisfies the Jordan
identity. In this project the product is written as \(xy\) in mathematical
prose and sometimes as \(x * y\) in code-oriented notes.
\end{definition}

\section{Indecomposables}

The guiding abelianization candidate is
\[
  Q(J) = J/J^2,
\]
where \(J^2\) denotes the linear span of products \(xy\). The slogan is useful
only after the ambient category and the unital convention have been fixed.

\begin{question}
What is the cleanest first setting for \(Q(J)=J/J^2\): absolute nonunital
Jordan algebras, augmented unital Jordan algebras, or a relative category?
\end{question}

\begin{note}
  In the absolute nonunital category, the abelian objects are precisely the
  Jordan algebras with zero multiplication. Indeed, a vector space \(V\) with
  product
  \[
      uv = 0
      \qquad
      \text{for all }u,v\in V
  \]
  is automatically a Jordan algebra. Therefore the natural abelianization of a
  nonunital Jordan algebra \(J\) is not a commutator quotient, but rather
  \[
      Q(J)=J/J^2.
  \]
  Here \(J^2\) is the linear span of all products. The quotient has zero
  multiplication because every product in \(J/J^2\) is represented by an element
  of \(J^2\).

  Let \(A\) be a unital Jordan algebra over \(k\). An \emph{augmentation} of \(A\)
is a unital Jordan algebra morphism
\[
    \varepsilon : A \longrightarrow k,
\]
where \(k\) is regarded as a unital Jordan algebra with its usual
multiplication.

The pair \((A,\varepsilon)\) is called an \emph{augmented unital Jordan
algebra}. Its \emph{augmentation ideal} is
\[
    A_+ := \ker(\varepsilon).
\]

Since \(\varepsilon(1_A)=1\), the canonical inclusion
\[
    k \longrightarrow A,
    \qquad
    \lambda \longmapsto \lambda 1_A
\]
splits \(\varepsilon\) as a map of \(k\)-vector spaces. Hence there is a vector
space decomposition
\[
    A \cong k\cdot 1_A \oplus A_+.
\]

Moreover, \(A_+\) is an ideal. Indeed, if \(a\in A_+\) and \(b\in A\), then
\[
    \varepsilon(ab)
    =
    \varepsilon(a)\varepsilon(b)
    =
    0,
\]
so \(ab\in A_+\). In particular, if \(a,b\in A_+\), then \(ab\in A_+\).
Thus \(A_+\) is naturally a nonunital Jordan algebra.

The augmentation ideal is important because, for a unital algebra \(A\), the
naive quotient
\[
    A/A^2
\]
is usually useless: since \(1_A\in A\), one typically has
\[
    A^2=A.
\]
Therefore, in the augmented unital setting, the correct analogue of
indecomposables is
\[
    Q(A)
    :=
    A_+/A_+^2,
\]
where
\[
    A_+ = \ker(\varepsilon).
\]
This is the same idea as in the nonunital setting, but applied to the
augmentation ideal rather than to the whole unital algebra.

Conversely, if \(J\) is a nonunital Jordan algebra, one may form its
unitalization
\[
    J^{+} := k\oplus J.
\]
The product is defined by
\[
    (\lambda,x)(\mu,y)
    :=
    \bigl(
        \lambda\mu,
        \lambda y+\mu x+xy
    \bigr),
\]
where \(\lambda,\mu\in k\) and \(x,y\in J\). The element
\[
    (1,0)
\]
is the unit. There is a natural augmentation
\[
    \varepsilon : J^{+}\longrightarrow k,
    \qquad
    \varepsilon(\lambda,x)=\lambda,
\]
and its kernel is
\[
    \ker(\varepsilon)=0\oplus J\cong J.
\]
Thus nonunital Jordan algebras may often be studied equivalently through
augmented unital Jordan algebras by passing between
\[
    J
    \qquad\text{and}\qquad
    k\oplus J.
\]
\end{note}

\section{Modules, Extensions, And Differentials}

The coefficient objects should be understood through square-zero extensions
\(J \ltimes M\), Beck modules over \(J\), and Jordan derivations. The main
research question is whether the Beck-module convention exactly matches the
classical Jordan-module convention used in the literature. (I need to ask my supervisor about this).

The expected Jordan analogue of Kahler differentials is a universal object
\(\Omega_{B/A}\) representing Jordan derivations. The corresponding cotangent
complex \(L^{\mathrm{Jord}}_{B/A}\) is still provisional until the resolution
framework is fixed.

\chapter{Workspace Architecture}

\section{Repository Map}

The repository separates mathematical claims, computational evidence, source
code, paper text, and private notes. This separation is part of the research
method: a computation may suggest a claim, but it does not by itself prove it.

\section{Theory And Claims}

\path{theory/} stores the evolving mathematical development. The directory
\path{theory/claims/} is the formal claim registry, and every theorem-like
statement should eventually be represented there with status, assumptions,
dependencies, proof sketch, references, and verification notes.

The current claim statuses are deliberately conservative. \path{CLAIM-0001}
and \path{CLAIM-0002} are draft claims, while \path{CLAIM-0003} through
\path{CLAIM-0005} are conjectural.

\section{Code, Tests, And Experiments}

\path{src/jordan_qh/} contains exact Python scaffolding for small examples,
linear algebra, identities, derivations, modules, omega candidates,
indecomposables, and toy Quillen checks. The matching tests live in
\path{tests/}.

\path{experiments/} records planned or executed computations. The experiment
numbers are stable workspace identifiers, not a logical dependency order.

\section{Paper, Formalization, And Private Notes}

\path{paper/} is the public-facing LaTeX skeleton and should only receive
statements that are aligned with checked or clearly labelled claim statuses.
\path{formal/} is a Lean scaffold, not part of the default verification path.

\path{note/} is a private learning area. This file may contain bilingual
explanations, partial understanding, and personal TODOs, but it should not
silently promote experiments or conjectures into proved results.

\chapter{Experiments}

\section{Reading Rule}

Each experiment should be read as computational evidence with an explicit
boundary. A result file may report a successful run, but the corresponding
mathematical claim still needs independent review.

\section{Experiment 001: Square-Zero Sanity Checks}

Experiment 001 has status \emph{run}. It checks one-dimensional split
square-zero extensions \(E = J \ltimes M\) with exact rational arithmetic.

The result is a sanity check for scalar actions in the idempotent and
zero-multiplication one-dimensional cases. It does not prove the Beck-module
classification and does not change the status of \path{CLAIM-0003}.

\section{Experiment 002: Two-Dimensional Jordan Algebras}

Experiment 002 has status \emph{not run yet}. Its goal is to enumerate or
import small two-dimensional Jordan algebra examples and test derivations,
indecomposables, and candidate differential modules.

\section{Experiment 003: Beck Module Identities}

Experiment 003 has status \emph{not run yet}. It is intended to symbolically
derive the identities imposed on \(M\) by the Jordan identity on a square-zero
extension \(J \ltimes M\).

\section{Experiment 004: TKK Comparison}

Experiment 004 has status \emph{not run yet}. Its purpose is to compare
intrinsic Jordan Quillen invariants with Lie homology of the associated TKK
Lie algebra on small examples.

\section{Experiment 005: Derived Indecomposables Toy Benchmark}

Experiment 005 has status \emph{run}. It benchmarks a one-generator toy model
for derived indecomposables on \(F(x)\), \(F(x)/(x^2-x)\), \(F(x)/(x^2)\),
\(F(x)/(x^3)\), and \(F(x)/(x^4)\).

The rule is that after applying \(Q\), only the linear part of a relation
survives. The current generated result says that all five planned examples
match their expected \(H_0\) and \(H_1\) dimensions.

\section{\texorpdfstring{Experiment 006: Presentation-Invariance Stress Test For \(J_3\)}{Experiment 006: Presentation-Invariance Stress Test For J3}}

Experiment 006 has status \emph{run}. It studies
\[
  J_3 = (t)/(t^3)
\]
through Presentation A, \(F(x)/(x^3)\), and a naive diagnostic for Presentation
B, \(F(x,y)/(y-x^2, xy, y^2)\).

The current result computes \(H_0=1, H_1=1\) for Presentation A. The naive
one-step diagnostic for Presentation B computes \(H_0=1, H_1=2\), exposing an
extra false class caused by the incomplete relation complex.

\section{\texorpdfstring{Experiment 007: Two-Step Cofibrant Replacement Toy For \(J_3\)}{Experiment 007: Two-Step Cofibrant Replacement Toy For J3}}

Experiment 007 has status \emph{run}. It keeps Presentation B from Experiment
006,
\[
  F(x,y)/(y-x^2, xy, y^2),
\]
and adds one explicit degree \(2\) generator \(b\) to kill the extra naive
class represented by \(a_3\).

The hand-coded two-step replacement uses degree \(0\) generators \(x,y\),
degree \(1\) generators \(a_1,a_2,a_3\), and differential
\[
\begin{aligned}
  d(a_1) &= y-x^2,\\
  d(a_2) &= xy,\\
  d(a_3) &= y^2,\\
  d(b) &= a_3 - x a_2 - (y+x^2)a_1 + x(xa_1).
\end{aligned}
\]
After applying indecomposables \(Q(-)=(-)/(-)^2\), the complex is
\[
  Q_2=\langle b\rangle \longrightarrow
  Q_1=\langle a_1,a_2,a_3\rangle \longrightarrow
  Q_0=\langle x,y\rangle
\]
with
\[
  d_Q(a_1)=y,\qquad d_Q(a_2)=0,\qquad d_Q(a_3)=0,\qquad d_Q(b)=a_3.
\]

The generated result records \(d_Q^2=0\), and the symbolic check behind the
degree \(2\) attachment reduces in this fixed one-generated example to
\[
  x^2x^2 - x(xx^2)=0,
\]
using power-associativity in the subalgebra generated by \(x\).

The computed homology dimensions are
\[
  H_0=1,\qquad H_1=1,\qquad H_2=0,
\]
with representatives \(x\) in \(H_0\), \(a_2\) in \(H_1\), and no \(H_2\)
representative. Compared with Experiment 006, the naive Presentation B
\(H_1\)-dimension drops from \(2\) to \(1\).

The boundary is essential: this verifies the fixed low-degree two-step repair
for \(J_3\) Presentation B through degree \(1\). It is not an automatic syzygy
search, not a general cofibrant replacement algorithm, and not a proof of full
presentation invariance.

\chapter{Experiment Summary}

\section{What The Current Runs Support}

The current executed experiments support the usefulness of exact arithmetic
small examples. Experiment 001 checks square-zero extension behavior in tiny
families, Experiment 005 validates a one-generator toy rule, and Experiment 006
shows why presentation issues cannot be ignored. Experiment 007 then repairs
the fixed \(J_3\) Presentation B diagnostic through degree \(1\) by adding one
explicit degree \(2\) generator.

\section{What The Current Runs Do Not Prove}

The successful toy runs do not prove the full Jordan Quillen homology theory.
They do not establish a general cofibrant replacement construction, a
presentation-invariance theorem, or an equivalence between intrinsic Jordan
Quillen homology and TKK Lie homology.

\section{Most Important Boundary}

The most important current warning is the pair of Experiments 006 and 007.
Experiment 006 shows that a naive one-step relation complex for Presentation B
has an extra false \(H_1\) class. Experiment 007 shows that this particular
false class can be killed by a hand-coded degree \(2\) generator, but the
repair is still example-specific and low-degree.

\section{Personal Takeaway}

The experiments are best used as guardrails for definitions. They are useful
when they expose impossible claims, missing hypotheses, or misleading
computational shortcuts. The 007 run is especially useful as a template for
what later automatic low-degree cell attachment would need to reproduce.

\chapter{Theory Roadmap}

\section{Abelianization}

The first theoretical target is to make \(Q(J)=J/J^2\) precise in the chosen
category. The related claim files are currently draft-level and should not be
quoted as checked theorems.

\section{Beck Modules And Square-Zero Extensions}

The next target is to decide whether Beck modules over a Jordan algebra match
the classical Jordan module identities. This is central because it controls the
coefficient category for cohomology.

\section{Derivations And Universal Differentials}

After the coefficient category is clear, the project should construct and test
universal Jordan differentials \(\Omega_{B/A}\). The universal property remains
a conjectural or proof-draft target until the module convention is settled.

\section{Cotangent Complex And Low Degrees}

The cotangent complex \(L^{\mathrm{Jord}}_{B/A}\) requires a resolution
framework, such as simplicial Jordan algebras, dg Jordan algebras over the
Jordan operad, or a comparison between both.

Experiment 007 suggests the next concrete computational target: replace naive
one-step relation complexes by a principled low-degree cell-attachment
procedure. The present evidence is only the fixed \(J_3\) two-step repair, so
it should guide implementation tests rather than be quoted as a theorem.

Low-degree cohomology should eventually explain \(D^0\) in terms of Jordan
derivations and \(D^1\) in terms of square-zero extensions. These statements
need separate claim files before they become paper-level assertions.

\section{Comparison Problems}

The comparison with Glassman cohomology and TKK Lie-theoretic constructions
should remain explicit and cautious. Agreement in small examples is useful,
but non-agreement may be equally informative.


\chapter{Current Repository Assessment}

\section{Snapshot Boundary}

This chapter records a working assessment of the repository state. It is a
private planning note, not a claim file and not proof evidence. Any statement
below that depends on generated results should be checked against the relevant
\path{results.json}, \path{experiment.md}, commit record, and verification
commands before it is quoted in public project documentation.

\begin{note}
The current project stage is best described as
\[
  \text{reproducible experimental scaffold}
  \quad+\quad
  \text{small exact computations}.
\]
It is not yet a completed theory of intrinsic Jordan Quillen homology. The
useful achievement is narrower but important: the repository now has enough
notation, claims, tests, and toy computations to expose false shortcuts and to
shape the next definitions.
\end{note}

\section{Overall Diagnosis}

The project goal is clear. The intended construction is intrinsic: use the
Jordan algebra category itself, rather than first replacing a Jordan algebra by
an associated Lie algebra such as the Tits--Kantor--Koecher algebra. The guiding
objects are derived abelianization, Jordan cotangent complexes, low-degree
Quillen cohomology, square-zero extensions, and comparisons with Glassman or
TKK-style constructions.

The repository discipline is also good. It asks for a strict distinction
between proved claims, proof drafts, conjectures, and computations. It also
requires exact arithmetic where possible and forbids fabricated experiment
output. This is especially important here because small Jordan examples can
look convincing even when a naive relation complex is not homotopically
invariant.

At the code level, the package is intentionally light: a Python package over
Python \(3.11\) or later, with no runtime dependencies and development tools
such as \path{pytest}, \path{ruff}, and \path{pre-commit}. The \path{Makefile}
exposes entries such as \path{test}, \path{lint}, \path{claims}, and
\path{experiments}. A current practical caveat is that \path{make experiments}
does not mean that every experiment is rerun: the top-level
\path{scripts/run_all_experiments.py} currently behaves as a conservative
status/reporting surface and prints that experiments are not run by default.

\section{Claim Status}

The theory layer remains deliberately conservative.

\begin{itemize}[leftmargin=*]
  \item \path{CLAIM-0001} is draft-level: zero-multiplication Jordan algebras
  are the expected abelian objects in the relevant nonunital setting.
  \item \path{CLAIM-0002} is draft-level: the indecomposables functor
  \(Q(J)=J/J^2\) is expected to be the abelianization or left adjoint in the
  chosen absolute nonunital setting.
  \item \path{CLAIM-0003} is conjectural: Beck modules over Jordan algebras
  should match a precise Jordan module convention, but the identity extraction
  still needs to be completed.
  \item \path{CLAIM-0004} is conjectural: universal Jordan differentials
  \(\Omega_{B/A}\) need a settled coefficient category and universal property.
  \item \path{CLAIM-0005} is conjectural: the transitivity triangle for the
  Jordan cotangent complex is not yet proved in the project.
\end{itemize}

\section{Experiment Status Matrix}

The current experiment landscape is:

\begin{description}[leftmargin=*,style=nextline]
  \item[001 square-zero, run.]
  A finite-grid sanity check for one-dimensional split extensions.
  \item[002 two-dimensional, not run.]
  Representative small examples still need to be computed.
  \item[003 Beck modules, not run.]
  Symbolic module identities are still a placeholder target.
  \item[004 TKK comparison, not run.]
  Comparison examples have not yet been computed.
  \item[005 derived \(Q\), run.]
  The one-generator toy benchmark matches the expected dimensions.
  \item[006 \(J_3\) stress test, run.]
  Naive Presentation B has an extra false \(H_1\) class.
  \item[007 two-step repair, run.]
  One hand-coded degree \(2\) generator kills that false class in the tested
  low-degree complex.
\end{description}

The status should be read from each experiment's own files, especially
\path{results.json} or \path{results.md}. The global runner is not yet a
complete experiment execution engine.

\section{Experiment 001: Split Square-Zero Extensions}

Experiment 001 studies one-dimensional split square-zero extensions
\[
  E = J\ltimes M,
  \qquad
  M^2=0.
\]
It tests two one-dimensional bases:
\[
  e^2=e
  \qquad\text{and}\qquad
  e^2=0,
\]
with \(M=km\) and scalar action
\[
  L_e(m)=\lambda m.
\]

The important methodological point is that the experiment uses exact rational
arithmetic and checks the Jordan identity on a finite grid of linear
combinations, not only on basis vectors. This catches a real failure mode:
\(\lambda=2\) can pass a basis-only smoke test but fail once general linear
combinations are tested.

For \(e^2=e\), the tested nonunital split extension passes for
\[
  \lambda\in\{0,\tfrac12,1\}.
\]
If \((e,0)\) is required to remain a unit in the extension, the surviving value
is \(\lambda=1\). For \(e^2=0\), the tested action is forced to
\[
  \lambda=0.
\]

This experiment supports the square-zero workflow behind \path{CLAIM-0003},
but it does not prove the Beck-module classification. Its main lesson is that
future module-identity checks must use symbolic derivation or a strong
polynomial identity check, not basis-vector testing alone.

\section{Experiment 005: One-Generator Derived Indecomposables}

Experiment 005 is the cleanest current Quillen-style toy computation. It
studies one-generator presentations
\[
  F(x),\quad
  F(x)/(x^2-x),\quad
  F(x)/(x^2),\quad
  F(x)/(x^3),\quad
  F(x)/(x^4).
\]

The toy rule is that after applying indecomposables \(Q\), only the linear part
of a relation enters the first relation complex. In other words, the current
toy complex is the two-term complex
\[
  k^{\{\text{relations}\}}
  \longrightarrow
  k^{\{\text{generators}\}},
\]
where the differential records relation linear terms. Exact row-rank
computations then give \(H_0\) and \(H_1\).

The current benchmark conclusions are:
\[
\begin{array}{c|c}
\text{presentation} & (H_0,H_1) \\
\hline
F(x) & (1,0) \\
F(x)/(x^2-x) & (0,0) \\
F(x)/(x^2) & (1,1) \\
F(x)/(x^3) & (1,1) \\
F(x)/(x^4) & (1,1).
\end{array}
\]

This gives a useful low-degree intuition for \(Q(J)=J/J^2\): a relation with a
linear term can kill \(H_0\), while a pure higher-order relation appears as a
naive \(H_1\) class in this one-step model. The boundary remains essential:
this is not a general cofibrant replacement engine and it does not establish
presentation invariance.

\section{\texorpdfstring{Experiment 006: Presentation Stress Test For \(J_3\)}{Experiment 006: Presentation Stress Test For J3}}

Experiment 006 compares two presentations of
\[
  J_3=(t)/(t^3).
\]
The first presentation is
\[
  \text{A: } F(x)/(x^3),
\]
and the second is
\[
  \text{B: } F(x,y)/(y-x^2,\;xy,\;y^2).
\]

Presentation A agrees with the one-generator toy computation and gives
\[
  H_0=1,\qquad H_1=1.
\]
The naive one-step diagnostic for Presentation B gives
\[
  H_0=1,\qquad H_1=2.
\]
The result is therefore not a failure of the project. It is a useful negative
test: the naive relation complex sees an extra false \(H_1\) class for
Presentation B. Without higher syzygy data or a genuine cofibrant replacement,
Presentation B cannot be used as evidence for presentation invariance.

\begin{note}
The practical target created by Experiment 006 is precise: identify the extra
class in the naive \(H_1\), then attach a higher cell or syzygy generator that
kills it without changing the intended algebra.
\end{note}

\section{\texorpdfstring{Experiment 007: Two-Step Toy Repair For \(J_3\)}{Experiment 007: Two-Step Toy Repair For J3}}

Experiment 007 performs exactly that fixed repair. It starts from Presentation
B and adds one degree \(2\) generator \(b\), designed to kill the redundant
class represented after applying \(Q\) by \(a_3\), the relation corresponding
to \(y^2=0\).

The low-degree indecomposable complex has the form
\[
  Q_2=\langle b\rangle
  \longrightarrow
  Q_1=\langle a_1,a_2,a_3\rangle
  \longrightarrow
  Q_0=\langle x,y\rangle,
\]
with
\[
  d_Q(a_1)=y,\qquad
  d_Q(a_2)=0,\qquad
  d_Q(a_3)=0,\qquad
  d_Q(b)=a_3.
\]
Hence \(b\) kills exactly the extra class that was visible in the naive
Presentation B computation.

The computed dimensions become
\[
  H_0=1,\qquad H_1=1,\qquad H_2=0.
\]
Thus the naive \(H_1\)-dimension \(2\) from Experiment 006 drops to \(1\), in
agreement with Presentation A through the tested degree.

This is the most important current mathematical evidence in the repository:
the project has moved from detecting a presentation artifact to hand-coding a
low-degree semi-free repair that removes it. The result must still be quoted
carefully. It is a fixed \(J_3\) computation, not an automatic syzygy search,
not a general cofibrant replacement algorithm, and not a proof of full
presentation invariance.

\section{Current Strengths}

The main strengths of the current repository are:

\begin{itemize}[leftmargin=*]
  \item The research boundary is explicit: computation is evidence, not proof.
  \item The claim files are conservative and do not overstate the theory.
  \item The package already has exact linear algebra utilities for small
  examples.
  \item The indecomposables code computes the rank of \(J^2\) and therefore
  \(\dim Q(J)=\dim J-\dim J^2\) in finite examples.
  \item The tests cover basic cases such as zero multiplication, a
  one-dimensional idempotent algebra, and a one-dimensional square-zero
  algebra.
  \item Experiments 006 and 007 give a concrete warning and repair pattern for
  presentation-sensitive toy computations.
\end{itemize}

\section{Current Risks}

The main risks are:

\begin{enumerate}[leftmargin=*]
  \item The coefficient theory is not yet settled. The files for modules,
  universal differentials, and TKK comparison remain draft or placeholder
  surfaces rather than completed theory.
  \item Experiment 007 is hand-coded. It should inspire the next algorithm, but
  it is not itself an automatic cell-attachment or syzygy-discovery procedure.
  \item Experiments 002, 003, and 004 are still the missing bridge between the
  one-generator toy computations and the broader Jordan algebra theory.
  \item The global experiment runner does not yet summarize the whole
  experiment state in a machine-checkable matrix.
\end{enumerate}

\section{Next Experimental Priorities}

\subsection*{Priority 1: Experiment 008}

The next best experiment is an automatic syzygy-discovery version of
Experiment 007. The goal is to input the three first-stage relations
\[
  a_1,\quad a_2,\quad a_3
\]
from Presentation B and automatically find the degree \(2\) attachment whose
image under \(Q\) kills \(a_3\). The expected output should reproduce
\[
  H_0=1,\qquad H_1=1,\qquad H_2=0,
\]
while explicitly saying that this is only a degree \(1/2\) toy replacement, not
a full cofibrant replacement.

\subsection*{Priority 2: Experiment 003}

The Beck-module identity experiment should be restored and made symbolic. The
starting point should be the Jordan identity
\[
  ((a^2b)a)=a^2(ba)
\]
inside a square-zero extension \(J\ltimes M\) with \(M^2=0\). The experiment
should extract the \(M\)-linear terms and compare the resulting operator
identities for \(L_x\) with the classical Jordan module axioms. This is central
for \path{CLAIM-0003}.

\subsection*{Priority 3: Experiment 002}

The two-dimensional example table should be completed in a modest,
representative way before attempting classification. Good first examples are:
the two-dimensional zero product algebra, \(ke\oplus km\) with \(e^2=e\), a
one-dimensional idempotent plus a zero ideal, and a small two-dimensional
nilpotent example. For each example, record \(\dim J^2\), \(\dim Q(J)\),
derivation dimensions, and candidate relations for \(\Omega\).

\subsection*{Priority 4: Experiment 004}

The TKK comparison should begin with two or three minimal examples and should
allow disagreement as a valid outcome. The repository rule is correct: one
should not identify intrinsic Jordan Quillen homology with TKK Lie homology
without either a proof or a counterexample.

\subsection*{Priority 5: Experiment Status Aggregator}

The script \path{scripts/run_all_experiments.py} should become a true
read-only status aggregator before it becomes an execution driver. A useful
first version would read each \path{experiment.md} and \path{results.json} or
\path{results.md}, then produce a matrix recording status, pass/fail data, and
the presentation-invariance flags that already appear in the JSON outputs. By
default it should avoid rewriting results.

\section{Working Conclusion}

The most important current mathematical progression is
\[
  \text{Experiment 006}
  \quad\longrightarrow\quad
  \text{Experiment 007}.
\]
Experiment 006 shows that a naive Presentation B relation complex for \(J_3\)
contains an extra false \(H_1\) class. Experiment 007 shows that a specific
two-step semi-free toy replacement kills that class and restores the expected
low-degree dimensions through degree \(1\).

The next phase should therefore not jump directly to paper-level theorems. It
should first turn the hand-coded success of Experiment 007 into a repeatable
low-degree method, while also completing the coefficient-theory and small
example experiments that make the theory meaningful.

\appendix
\chapter{Notation And Literature}

\section{Notation Anchor}

The canonical notation table is \path{theory/latex_notation.md}. This private
note should follow that file instead of inventing competing notation.

\section{Basic Notation}

The most important current notations are: ground field \(k\), Jordan algebra
\(J\), module \(M\), product span \(J^2\), indecomposables \(Q(J)=J/J^2\),
cotangent complex \(L^{\mathrm{Jord}}_{B/A}\), homology
\(D_n^{\mathrm{Jord}}(B/A,M)\), and cohomology
\(D^n_{\mathrm{Jord}}(B/A,M)\).

\section{Literature Notes}

The curated literature notes live in \path{literature/notes/}. They are the
Codex-facing layer for references; local PDFs remain under
\path{literature/papers-local/} and should be checked by the human reader when
precise formulas or citations matter.

\section{Verification Reminders}

For Python changes, the default checks are:
\begin{itemize}
  \item \path{python -m pytest}
  \item \path{python -m ruff check .}
  \item \path{python scripts/check_claims.py}
\end{itemize}
For this private note, the relevant check is that the LaTeX file compiles.


\end{document}
